\subsubsection*{Quantization steps and the \(W/Z\) landing}

The table below names the Standard-Model field-theory steps used in this
section, together with what each one implies, what its construction requires,
and its present status. The labels \(\mathrm{QFT}\text{-}\mathrm{Q0}\) through
\(\mathrm{QFT}\text{-}\mathrm{Q4}\) are local to this section and stay distinct
from the particle-receipt clauses \((Q1)\)--\((Q3)\). Their dependency
structure is a graph rather than a single ordered ladder:
\begingroup\small
\[
\begin{array}{rcl}
\mathsf{declared\ finite\ action}&\longrightarrow&\mathsf{QFT\!-\!Q1}
 \longrightarrow
 \left\{\begin{array}{l}
 \mathsf{QFT\!-\!Q2E/QFT\!-\!Q2H},\\
 \mathsf{QFT\!-\!Q3\ BV/ST}\longrightarrow
 \mathsf{strict\ finite\!-\!order\ }W/Z;
 \end{array}\right.\\[3pt]
\mathsf{nonperturbative\ observable\ tower}&\longrightarrow&
\mathsf{QFT\!-\!Q4\ OS}\\
&&\longrightarrow\mathsf{QFT\!-\!Q4\ resonance}.
\end{array}
\]
\endgroup
Native provenance on QFT-Q1 additionally requires QFT-Q0 and a
source/action-identity receipt.  No QFT-Q2-to-QFT-Q3,
QFT-Q3-to-QFT-Q2, or QFT-Q3-\(W/Z\)-to-QFT-Q4 arrow is automatic.

\begin{center}
\small
\begin{tabular}{@{}p{0.10\textwidth}p{0.31\textwidth}p{0.25\textwidth}p{0.24\textwidth}@{}}
\toprule
Tier & Mathematical implication & Required construction & Present status \\
\midrule
QFT-Q0 & finite representation, charge, anomaly, lattice, and selector consequences on the declared packet
& physical source selection and attachment
& conditional finite pass; producer incomplete \\
QFT-Q1 & finite local classical \(G_6\) action is gauge invariant and local; the standard tree kernel also needs canonical continuum normalization
& source-selected action, coefficients, regulator, normalization, and ancestry
& implication specified; producer open \\
QFT-Q2-E/H & an equivariant determinant-line section or a noncollapsing constrained Hamiltonian supplies an exact finite quantum object
& full operator, measure/current or Hamiltonian/nonvacuum packet, and refinement controls
& criteria specified; constructions open \\
QFT-Q3 & stable anomaly-free BV/ST theory is formally restorable order by order; strict \(W/Z\) poles land here
& counterterm basis, matching and Faddeev--Jackiw (FJ) engines, identities, currents, and numerical freeze
& implication specified; imported validation possible; OPH producer open \\
QFT-Q4 & OS reconstruction and a separate continued-sheet resonance theorem
& reflection-positive tower and analytic-continuation packet
& implications specified; constructions open \\
\bottomrule
\end{tabular}
\end{center}

\begin{theorem}[Finite local classical \(G_6\) action at QFT-Q1]
\label{thm:qft-q1-finite-local-g6}
Let \(K_r\) be a finite oriented spin four-complex with bounded incidence,
positive cell weights, declared boundary conditions, paired edge orientations,
and declared spin transports. Put \(U_e\in
G_6=S(U(3)\times U(2))\) on oriented edges and the declared Higgs and
left-handed matter variables on vertices. Require every matter action to
descend to a well-defined representation of \(G_6\), every plaquette term to be a
declared class function of its \(G_6\) holonomy, and every finite difference
to be gauge covariant with defined endpoint data. Then a finite sum of these
plaquette, Higgs, fermion, and invariant Yukawa terms is local and exactly
gauge invariant.
\end{theorem}

\begin{proof}
Plaquette holonomies transform by conjugation, covariant edge differences
transform at their endpoints, and class functions and invariant contractions
remove those transformations. In integer hypercharge normalization the three
Yukawa sums are
\[
1+3-4=0,\qquad 1-3+2=0,\qquad -3-3+6=0.
\]
Every term has bounded cell support.
\end{proof}

\begin{corollary}[Canonically normalized electroweak tree kernel]
\label{cor:qft-q1-ew-tree-kernel}
If, in addition, the long-wavelength map sends the finite Higgs kinetic form
to \((D_\mu H)^\dagger D^\mu H\) with canonical generator normalization and
broken background \(H_0=2^{-1/2}(0,v)^T\), then
\[
w=\frac{g^2v^2}{4},\qquad
\mathcal M_N^2=\frac{v^2}{4}
\begin{pmatrix}g^2&-gg'\\-gg'&g'^2\end{pmatrix},
\qquad z=\frac{(g^2+g'^2)v^2}{4},
\]
and the other neutral eigenvalue is zero.
\end{corollary}

\begin{remark}[QFT-Q1 boundary]
This is a classical existence template after the complex, fields,
coefficients, boundary data, and normalization bridge are supplied. It does
not show that OPH selects them, construct a chiral measure, remove mirrors or
doublers, or produce a complex pole.
\end{remark}

\begin{theorem}[QFT-Q2-E equivariant determinant-line criterion]
\label{thm:qft-q2e-determinant-line}
At a fixed finite stage, let \(D_r(U)\) be a local gauge-covariant,
\(\gamma_5\)-Hermitian Ginsparg--Wilson operator on a connected admissible
gauge-field component on which the chiral-projector rank is constant, with
the declared spectral gap and locality bounds. A basis-independent,
gauge-invariant finite chiral measure exists precisely when the determinant
line admits a nowhere-zero gauge-equivariant section in the required locality
and smoothness class. For a nonfree gauge action this is an equivariant
statement over the action groupoid, including stabilizer actions. In local
connection form the measure current must reproduce the projector curvature
and obey global loop integrability. Flatness with trivial holonomy is only a
sufficient special case.
\end{theorem}

\begin{proof}
Changes of Weyl basis are transition functions of the determinant line. A
nowhere-zero equivariant section cancels them and descends to the gauge
quotient; conversely a basis-independent gauge-invariant phase supplies
compatible nonzero fiber vectors. The curvature equation and loop condition
are the local and global integrability conditions for that section.
\end{proof}

\begin{theorem}[QFT-Q2-H finite Hamiltonian soundness]
\label{thm:qft-q2h-hamiltonian-soundness}
Let \(\mathcal H_{\rm kin}\) be finite dimensional and let local Gauss
generators exponentiate to the complete local \(G_6\) action. Define
\[
C_G=\sum_{x,a}(G_x^a)^\dagger G_x^a,\qquad
P_{\rm phys}=\mathbf1_{\{0\}}(C_G),
\qquad
1<\operatorname{rank}P_{\rm phys}<\dim\mathcal H_{\rm kin}.
\]
Suppose a bounded-range self-adjoint \(H_r\) commutes with this action and has
a unique certified physical ground state \(\Omega_r\) with a positive gap.
If a bounded self-adjoint gauge-invariant observable \(O_r\) has strictly
positive variance in \(\Omega_r\), and the packet also supplies its claimed
chiral index, positive mirror gap, primitive completeness, and
complement-complete refinement controls, then
\(P_{\rm phys}\mathcal H_{\rm kin}\) is a nonvacuous finite unitary gauge
theory with a positive-energy nonvacuum physical excitation and the declared
chiral/mirror properties.
\end{theorem}

\begin{proof}
The vector
\[
(O_r-\langle\Omega_r,O_r\Omega_r\rangle)\Omega_r
\]
is physical, orthogonal to the ground state, and nonzero by the variance
condition. The remaining conclusions are exactly the separately supplied
index, gap, completeness, and refinement clauses. Thus neither an identity
projector nor a vacuum-only physical sector passes.
\end{proof}

\begin{remark}[QFT-Q2 boundary]
The QFT-Q2-E and QFT-Q2-H results are criteria and soundness theorems. The
current OPH corpus supplies neither full-\(G_6\) construction. Anomaly
arithmetic alone instantiates neither theorem.
\end{remark}

\begin{theorem}[Formal QFT-Q3 Slavnov--Taylor restoration]
\label{thm:qft-q3-formal-st-restoration}
Let \(S_0\) be the complete canonically normalized Standard-Model BV action,
including the gauge-fixing and antifield sectors, and suppose
\((S_0,S_0)=0\). Fix a regulator/scheme satisfying the quantum action
principle, locality, and the declared power counting. Assume stability under
renormalization in a complete permitted counterterm basis with fixed
normalization conditions, classification of the applicable local
ghost-number-one BRST cohomology, vanishing of the declared perturbative
anomaly class, and a separate check of all applicable global anomalies. Then
finite local counterterms may be chosen recursively so that the formal series
\[
\Gamma=S_0+\sum_{n\ge1}\kappa^n\Gamma_n,\qquad
\kappa=(16\pi^2)^{-1},
\]
satisfies the renormalized Slavnov--Taylor identity order by order.
\end{theorem}

\begin{proof}
At order \(n\), the quantum action principle makes the breaking a local
ghost-number-one functional \(\Delta_n\). Wess--Zumino consistency makes it
closed under the linearized Slavnov--Taylor operator. The cohomology
classification splits it into an anomaly representative plus an exact term.
Anomaly clearance removes the first, and an allowed finite counterterm
cancels the second. Stability and the normalization conditions fix the
remaining invariant freedom, so induction proves the formal statement.
\end{proof}

\begin{remark}[QFT-Q3 boundary]
This is a formal power-series theorem, not a convergence theorem, QFT-Q2
construction, or QFT-Q4 Wightman construction. A numerical implementation
must produce its regulator-specific restoration transcript and verify
the Ward, Slavnov--Taylor, and Nielsen identities.
\end{remark}

\begin{lemma}[Conditional first-order FJ coordinate change]
\label{lem:qft-q3-fj-reparametrization}
Freeze the potential normalization, bare VEV shift, tadpole prescription,
field and parameter counterterms, mass arguments, mixing coordinates, and
gauge-fixing convention. If the complete finite change is
\[
p_L=p_F+\kappa\,\delta p^{(1)}+O(\kappa^2)
\]
and \(s(p)=s_0(p)+\kappa s_1(p)+O(\kappa^2)\), equality of the exact pole in
the two coordinates implies
\[
s_{1,F}=s_{1,L}+\delta p^{a(1)}\partial_as_0.
\]
\end{lemma}

\begin{proof}
Substitute \(p_L(p_F)\) and Taylor expand through first order. The sum must
include every transformed parameter, normalization, mass argument,
counterterm, and mixing coordinate; a VEV-only substitution is insufficient.
\end{proof}

\begin{theorem}[Strict charged and neutral pole coefficients]
\label{thm:qft-q3-strict-wz-coefficients}
On one scheme, contribution mask, and resonance sheet, write
\[
\Gamma_W^T=s-w+\kappa\Pi_{WW}^{(1)}
+\kappa^2\Pi_{WW}^{(2)}+O(\kappa^3).
\]
For \(s_W=w+\kappa s_{W,1}+\kappa^2s_{W,2}+O(\kappa^3)\),
\[
s_{W,1}=-\Pi_{WW}^{(1)}(w),\qquad
s_{W,2}=\Pi_{WW}^{(1)}(w)\Pi_{WW}^{(1)\prime}(w)
-\Pi_{WW}^{(2)}(w).
\]
For the massive root of the full photon--\(Z\) matrix with tree value
\(z\ne0\),
\[
s_{Z,1}=-\Pi_{ZZ}^{(1)}(z),
\]
\[
s_{Z,2}=\Pi_{ZZ}^{(1)}(z)\Pi_{ZZ}^{(1)\prime}(z)
-\Pi_{ZZ}^{(2)}(z)
+\frac{\Pi_{ZA}^{(1)}(z)\Pi_{AZ}^{(1)}(z)}{z}.
\]
Thus the one-loop-squared neutral mixing product is excluded at strict one
loop and is one mandatory term of the complete strict-two-loop mask.
\end{theorem}

\begin{proof}
Insert the root series and compare powers of \(\kappa\). In the neutral
sector use the Schur complement of the photon block; both off-diagonal
entries begin at order \(\kappa\).
\end{proof}

\begin{theorem}[Nielsen control and physical current pole at QFT-Q3]
\label{thm:qft-q3-nielsen-current-pole}
Suppose the inverse matrix and Nielsen insertions are holomorphic near a
simple massive root on the frozen sheet and, through retained order \(N\),
\[
\partial_\eta\Gamma^T
=\Lambda_\eta\Gamma^T+\Gamma^T\widetilde\Lambda_\eta
+O(\kappa^{N+1}).
\]
Then
\[
\partial_\eta\det\Gamma^T
=\operatorname{tr}(\Lambda_\eta+\widetilde\Lambda_\eta)
\det\Gamma^T+O(\kappa^{N+1}),
\qquad
\partial_\eta s_p=O(\kappa^{N+1}).
\]
If the simple left/right kernel vectors \(\ell,r\) have
\(\ell^\dagger\Gamma^{T\prime}(s_p)r\ne0\), and the dressed renormalized
BRST-invariant current vertices obey
\(J_L^\dagger r\ne0\) and \(\ell^\dagger J_R\ne0\), then the
gauge-invariant current amplitude has Laurent residue
\[
\frac{(J_L^\dagger r)(\ell^\dagger J_R)}
{\ell^\dagger\Gamma^{T\prime}(s_p)r}.
\]
\end{theorem}

\begin{proof}
Use
\(\partial_\eta\det\Gamma
=\operatorname{tr}(\operatorname{adj}\Gamma\,\partial_\eta\Gamma)\);
the adjugate identity remains valid at a singular matrix and avoids dividing
by the determinant. The simple-root implicit equation gives the
omitted-order gauge variation. The rank-one Laurent expansion of
\(\Gamma^{-1}\), contracted with the dressed current vertices, gives the
amplitude pole. This is not a positivity claim for an unstable elementary
field.
\end{proof}

\begin{theorem}[Gauge-invariant QFT-Q4 reconstruction]
\label{thm:qft-q4-gauge-invariant-os}
A compatible cofinal Schwinger family for a complete declared
gauge-invariant observable algebra that satisfies distributional convergence,
Euclidean covariance, graded symmetry/locality, reflection positivity,
clustering, growth/regularity, noncollapse, and refinement Cauchy control
reconstructs a positive Hilbert space, cyclic vacuum, positive-energy
translations, and the corresponding observable-sector Wightman distributions
and local graded net.
\end{theorem}

\begin{remark}
This observable-sector implication does not by itself construct colored local
fields, charged infrared sectors, confinement, asymptotic completeness, an
\(S\)-matrix, or a second-sheet resonance.
\end{remark}

\begin{theorem}[QFT-Q4 resonance and residue stability]
\label{thm:qft-q4-resonance-stability}
On one common continued sheet, write
\[
G_r(s)=\frac{N_r(s)}{D_r(s)}+G_{r,\rm reg}(s).
\]
Let a Jordan contour and its interior lie in a common domain on which
\(D_r,N_r,G_{r,\rm reg}\) and their limits are holomorphic. Assume locally
uniform convergence of these data and \(D_r'\), one simple enclosed zero
\(s_r\), uniform nonzero contour and derivative bounds, and the Rouch\'e
inequality
\[
\sup_C|D-D_r|<\inf_C|D_r|.
\]
Then the limiting denominator has one simple zero \(s_*\), \(s_r\to s_*\),
and, if \(N(s_*)\ne0\),
\[
\operatorname*{Res}_{s=s_r}G_r(s)
=\frac{N_r(s_r)}{D_r'(s_r)}
\longrightarrow
\frac{N(s_*)}{D'(s_*)}\ne0.
\]
\end{theorem}

\begin{proof}
Rouch\'e preserves the zero count for the holomorphic denominators.
Compactness and uniqueness give root convergence. Uniform numerator and
derivative convergence with the lower derivative bound gives residue
convergence. Ordinary uniform convergence of the meromorphic quotient through
its own poles is neither assumed nor valid.
\end{proof}

\begin{remark}[Producer boundary]
These results specify conditional implications only. The current corpus does
not instantiate the source-selected QFT-Q1 action, either full QFT-Q2 object,
the QFT-Q3 matching/FJ/two-engine/current and uncertainty packet, or the
QFT-Q4 tower and continued-sheet data.
\end{remark}
